This work brings us closer to applying computer-assisted, internal,
type-theoretic reasoning to a category,
such as in the category of cubical sets,
the category of groupoids, and the category of categories.
Roughly speaking, we would like to
provide efficient \emph{syntactic} proofs about \emph{semantic} objects
by writing a proof in a {domain-specific language},
and have that proof type-checked and translated
into a model of our choosing (which is also formalised in a theorem prover).
The steps we make towards this general goal are both in 
furthering the type theoretic analysis of certain categories
(\Cref{sec:exponentiable,sec:polynomials,sec:path-types,sec:cat-as-types}),
as well as implementing computer-assisted
syntax-semantic reasoning (\Cref{sec:hottlean}).

\subsection*{Background}

A \emph{domain specific language} (DSL) in our context
will mean a computer language implementing a type theory that is
tailored towards reasoning about specific models of interest.
By using a domain-specific language,
we can reason about models in a way that not only allows us to identify
when the same construction appears in different contexts,
but also to hide away details that are irrelevant to the construction being made,
elucidating the essence of a proof.

\emph{Martin-L\"of type theory} (MLTT) was developed
as a new intuitionistic foundation for mathematics \cite{martin1975,martin1998}.
It admits various extensions,
including Extensional Type Theory and Homotopy Type Theory (HoTT),
and has models in presheaves \cite{hofmann1997, awodey2024},
groupoids \cite{hofmann1995},
simplicial sets \cite{kapulkin2021} and cubical sets
\cite{bezem2014, cohen2015, awodey2026cartesian}.
One component of the HoTTLean project \cite{hua2025}
is the design of a domain-specific language implementing MLTT.

A model of MLTT necessarily admits certain categorical structure:
the presence of $\Sigma$-types in the type theory 
corresponds to a compositional structure on type families;
$\Pi$-types correspond to right adjoints of pullback functors along type families;
$\Id$-types correspond to functorial factorisation and lifting conditions.
However, as a target for interpreting the syntax of type theory,
these categorical descriptions are only stable under substitution (or pullback)
up-to-isomorphism.
This is known as the \emph{coherence problem}.

One solution to the coherence problem is to
provide and check explicit coherence equations.
This is the approach of \emph{categories with families} (CwF) \cite{clairambault2011}.
\emph{Natural models} \cite{awodey2018}
offer a simplification of CwFs
by capturing each type former as a \emph{universal} or \emph{generic}
construction on a classifier of type families,
by employing the powerful machinery of polynomial functors.

We will refer to the natural models style of semantics as \emph{algebraic}
(as in \emph{algebraic type theory} \cite{awodey2025}),
and use the word \emph{unalgebraic} in contrast.
This thesis will make the case for using both unalgebraic and algebraic semantics
(mainly for the sake of formalisation in a proof assistant):
an unalgebraic model is well-suited for writing an interpretation function
from the syntax to the model,
whereas an algebraic model is easier to construct.
Translating between the two is straightforward,
making this approach fairly practical.

In a natural model,
the classifier of all type families is necessarily \emph{external},
meaning that the object that classifies types lives in the category of presheaves
on contexts, rather than the category of contexts itself.
Unlike natural models, our approach will be \emph{internal},
meaning that we replace universal constructions on an external classifier
of all type families (in presheaves) with universal constructions
on classifiers of small type families (in the category itself).
Type theoretically,
this means that all our types are classified by a \emph{universe}.
This creates certain complications:
the category we work in is not necessarily locally Cartesian closed;
it may not even have all finite limits.
Fortunately, we demonstrate in \Cref{sec:exponentiable,sec:polynomials}
that these complications can be dealt with,
and in a way that allows us to extend our type theoretic analysis
to the category of categories $\Cat$ and the category of topological spaces $\Top$.

% For example, Lawvere algebraic theories are used to reason about
% equationally axiomatised mathematical objects such as 
% groups, rings, and modules \cite{lawvere1963}.
% Toposes admit internal reasoning in geometric logic \cite{maclane2012}.
% Simply typed $\lambda$-calculus is an internal language for
% Cartesian closed categories \cite{lambek1988}.

\subsection*{Overview of the thesis}

A significant portion of this work was motivated by
developments in HoTTLean project \cite{hua2025}.
At the very surface,
HoTTLean is a proof assistant, 
implemented as a domain-specific language (DSL) in Lean
using Lean's extensible metaprogramming features.
This DSL uses a deep embedding of MLTT in Lean,
which is then interpreted into a generic model of MLTT.
\Cref{sec:hottlean} details our contribution to this project
in formalising both the general semantics of MLTT,
as well as the Hofmann-Streicher groupoid model \cite{hofmann1995}
as a particular target for the interpretation.
This ongoing project aims to eventually fully implement
the idea of using MLTT for internal language reasoning,
with automation handling the burden of translating internal language proofs
into the model.

One outcome of HoTTLean was Steve Awodey's
cubically-inspired idea of using ``path types''
to simplify the construction of identity types in the groupoid model
(and other models as well) \cite{awodey2026}.
Instead of constructing identity types by hand,
it is enough to have path types
(which can be deduced from a closure condition on fibrations),
and a ``Hurewicz'' path-lifting condition on fibrations.
These conditions are often fairly easy to check in a model.
As a result,
identity types in the groupoid model were soon fully
formalised using path types in HoTTLean.
The path types conditions can be rephrased in a style
reminiscent of modal type theory \cite{licata2016,licata2017},
which we present in \Cref{sec:path-types}.
The main theorem of \Cref{sec:path-types} follows.
\begin{theorem*}[\labelcref{thm:id-elim}]
  Suppose $(\CC,\RR)$ is a $\pi$-preclan (\labelcref{def:pi-preclan}).
  If $\uu : \Ulow \to \U$ is an $\RR$-algebraic universe (\labelcref{def:algebraic-universe}) that
  admits $\Path$-types (\labelcref{def:mltt-alg-path-types-2})
  and a normal Hurewicz structure (\labelcref{def:hurewicz-defs}),
  then the universe also admits $\Id$-types (\labelcref{def:algebraic-id-types}).
\end{theorem*}
 
After various experiments in HoTTLean,
it became apparent that a model of MLTT (with a hierarchy of universes)
could be formulated in a manner that was internal
to the category of contexts $\CC$ itself \cite{hua2026},
rather than in the category of presheaves on contexts,
as is more typical (see \cite{clairambault2011, awodey2018}).
Roughly, this meant that we replaced the presheaf classifying all types
$\Ty : \CC^\op \to \Set$
with a hierarchy of representable presheaves $\Ty_0 \to \Ty_1 \to \Ty_2 \to \dots$,
or better yet, just with the objects (contexts)
$U_0 \to U_1 \to U_2 \to \dots$ that represent them.
This approach,
in part inspired by Uemura's notion of a representable map category
\cite{uemura2023},
proved to significantly simplify the construction of a model,
but necessitated a more general theory of polynomial functors,
since this simplification relied on an \emph{algebraic} construction
(in the sense of \cite{awodey2025}, as explained above).

Following a doctrine of working both \emph{internally} and \emph{algebraically}
has naturally led to various other results.
This includes a more general theory of exponentiability
with respect to a class of maps (a ``preclan''),
presented in \Cref{sec:exponentiable},
as well as the ensuing theory of polynomial functors,
presented in \Cref{sec:polynomials}.
Reid Barton was quick to point out a host of examples
in which the theory of exponentiability could be applied.
Insights from those examples paved the way for the rest of
the theory's development.
The main theorem, providing two equivalent definitions of
a map being \emph{$\RR$-exponentiabile} is given below.

\begin{theorem*}[\labelcref{theorem:exponentiable}]
  Let $(\CC,\RR)$ be a preclan (\labelcref{def:preclan})
  and $f : Y \to X$ a morphism in $\CC$.
  The following are equivalent.
  \begin{enumerate}
  \item
    The Beck--Chevalley condition (\labelcref{def:beck-chevalley}) holds at every pullback of
    $f$ in the preclan $(\CC,\RR)$
    (whenever such pullbacks exist).
  \item
    The pushforward in presheaves $\Pi_f : (\wCC/Y, \wRR) \to (\wCC/X, \wRR)$
    is a preclan morphism (\labelcref{def:preclan}).
  \end{enumerate}
\end{theorem*}

A new kind of dependent type theory emerges out of
the considerations of \Cref{sec:exponentiable},
which involves two different kinds of types and
which differs from MLTT in the behaviour of $\Pi$-types.
A fairly well-studied example of this is in the category of categories $\Cat$.
In short, to form a $\Pi$-type $\Pi_A B$ in $\Cat$,
$B$ must be a covariant type if $A$ is a contravariant type,
resulting in a covariant type $\Pi_A B$ (and dually).
In $\Cat$, $\Sigma$-types work in the same way as usual.
However, MLTT $\Id$-types (intensional identity types) are replaced with
$\Hom$-types that have two elimination rules,
originally due to North \cite{north2019},
relating to twisted arrow categories and various algebraic weak factorisation systems on $\Cat$.
We develop this new kind of algebraic model in \Cref{sec:cat-as-types}.
We summarise this model in the following statement.

\begin{theorem*}[\labelcref{ex:ACM-of-opfibrations}]
  There is an algebraic categorical model (\labelcref{def:ACM}) based in $\Cat$
  with a universe (\labelcref{def:ACM-universe}) classifying split opfibrations with
  $\Unit$-types, $\Sigma$-types, $\Pi$-types,
  and $\Hom$-types (\labelcref{def:ACM-universe-sigma,def:ACM-universe-pi,def:ACM-universe-hom}).
  Dually, the same holds for the universe classifying split fibrations.
\end{theorem*}

\subsection*{Structure of the thesis}
\Cref{sec:exponentiable} presents exponentiability conditions with
respect to a class of maps.
\Cref{sec:polynomials} develops polynomial functors,
using the definitions from \Cref{sec:exponentiable}.
\Cref{sec:hottlean} summarises the formalisation of the groupoid model in the HoTTLean project,
and introduces unalgebraic and algebraic models of MLTT.
The formalisation uses the theory of polynomials in $\pi$-clans,
which is an application of \Cref{sec:polynomials}.
\Cref{sec:path-types} presents an alternative construction of
algebraic identity types (from \Cref{sec:hottlean}) via path types.
Finally, \Cref{sec:cat-as-types} develops a new kind of algebraic model
in the category of categories,
providing a case study of an algebraic model based on
a $\tau$-clan (from \Cref{sec:exponentiable}).